% ============================================================
%  06_format_sortie.tex — Section 6 : Format de sortie et resultats
% ============================================================
\section{Format de sortie et résultats}

\subsection{Structure d'une solution}

Chaque solution est un dictionnaire \texttt{\{hexagone\_id: taille\}} où la taille est 5
(pentagone), 6 (hexagone inchangé) ou 7 (heptagone). Par exemple, pour \texttt{first.graph}
(3 hexagones) :

\begin{lstlisting}[style=benzai]
solution 1: v0=5 v1=6 v2=7
solution 2: v0=7 v1=6 v2=5
\end{lstlisting}

Cela signifie : l'hexagone 0 devient un pentagone, l'hexagone 1 reste un hexagone,
l'hexagone 2 devient un heptagone (et vice versa pour la solution 2).

\subsection{Résultats expérimentaux}

\subsubsection{\texttt{first.graph} — 3 hexagones linéaires}

\begin{center}
\begin{tabular}{lr}
  \toprule
  Caractéristique & Valeur \\
  \midrule
  Hexagones & 3 \\
  Carbones & 14 \\
  Arêtes du dual & 2 \\
  Hexagones gelés & 0 \\
  Hexagones libres & 3 \\
  Générateurs Aut($\GD$) & 1 (réflexion) \\
  \textbf{Solutions distinctes} & \textbf{2} \\
  \bottomrule
\end{tabular}
\end{center}

Les 2 solutions : $v_0 = 5, v_1 = 6, v_2 = 7$ et $v_0 = 7, v_1 = 6, v_2 = 5$.
La contrainte C1 impose $n_5 = n_7$, donc le central reste à 6.

\subsubsection{\texttt{second.graph} — 16 hexagones}

\begin{center}
\begin{tabular}{lr}
  \toprule
  Caractéristique & Valeur \\
  \midrule
  Hexagones & 16 \\
  Carbones & 60 \\
  Arêtes du dual & 19 \\
  Hexagones gelés & 5 \\
  Hexagones libres & 11 \\
  Générateurs Aut($\GD$) & 1 \\
  \textbf{Solutions distinctes} & \textbf{41} \\
  \bottomrule
\end{tabular}
\end{center}

Les 41 solutions ont été énumérées en moins d'une seconde par ACE. La répartition des
tailles sur l'ensemble des solutions montre que la majorité des hexagones restent à 6, avec
des substitutions ponctuelles en pentagones et heptagones.
